\chapter{Monte-Carlo}\label{ch:qm:monte-carlo}

La saisie de nuit valorise 40\,000 transactions à trajectoires dépendantes avec 100\,000 trajectoires chacune. L'\omterm{def:qm:estimation:estimator}{erreur type} sur le livre est de 80\,000 dollars, et le responsable des risques veut la voir descendre sous 20\,000 d'ici demain.
L'erreur de Monte-Carlo décroît comme l'inverse de la racine carrée du nombre de trajectoires~; diviser l'erreur par quatre exige donc de multiplier le nombre de trajectoires par seize, et seize fois plus de machines ne sont pas disponibles. La
solution consiste à rendre chaque trajectoire plus utile. Ce chapitre construit l'\omterm{def:qm:estimation:estimator}{estimateur} et ses générateurs, puis les techniques de \omterm{def:qm:monte-carlo:mc}{réduction de variance} qui modifient la constante (\omterm{def:qm:monte-carlo:vr}{variables antithétiques},
\omterm{def:qm:monte-carlo:vr}{variables de contrôle}, stratification, \omterm{def:qm:monte-carlo:is}{échantillonnage préférentiel}), les points à discrépance faible qui modifient le taux, et les méthodes de discrétisation et multiniveau qui déterminent le coût d'une trajectoire.

\section{Estimateurs et générateurs}

\begin{definition}[Méthode de Monte-Carlo, réduction de variance]\label{def:qm:monte-carlo:mc}
La \emph{méthode de Monte-Carlo}\index{méthode de Monte-Carlo} estime $\mu = \E[f(X)]$ par la moyenne $\hat\mu_n = \frac1n\sum_{i=1}^nf(X_i)$ de tirages indépendants. La \emph{réduction de
variance}\index{réduction de variance} désigne tout changement d'\omterm{def:qm:estimation:estimator}{estimateur} qui conserve $\E[\hat\mu_n] = \mu$ (ou rend le biais négligeable) et abaisse sa variance à coût égal.
\end{definition}

\begin{proposition}[Erreur et coût de l'estimateur de Monte-Carlo]\label{prop:qm:monte-carlo:clt}
Si $\sigma^2 = \Var(f(X)) < \infty$, alors $\sqrt n(\hat\mu_n - \mu) \Rightarrow \mathcal N(0, \sigma^2)$, de sorte que $\hat\mu_n \pm 1.96\,s_n/\sqrt n$ est un intervalle asymptotique à 95~\%. Si un tirage coûte $c$, le
coût pour atteindre une \omterm{def:qm:estimation:estimator}{erreur type} $\epsilon$ est $c\sigma^2/\epsilon^2$~: parmi deux \omterm{def:qm:estimation:estimator}{estimateurs sans biais}, le meilleur est celui dont le produit de la variance par le coût par tirage est le plus petit.
\end{proposition}

\begin{proof}
La première affirmation est le théorème central limite pour des tirages indépendants et identiquement distribués, combiné au lemme de Slutsky pour l'\omterm{def:qm:estimation:estimator}{estimateur} $s_n$. La seconde~: $\sigma/\sqrt n = \epsilon$ exige $n =
\sigma^2/\epsilon^2$ tirages.
\end{proof}

Le taux $n^{-1/2}$ ne dépend pas de la dimension, ce qui permet à la \omterm{def:qm:monte-carlo:mc}{méthode de Monte-Carlo} de valoriser des moyennes sur 64 dates et des paniers de 30 actifs qu'aucune grille ne peut atteindre~; la constante $\sigma^2c$ est l'endroit où se
concentre le travail. Sur la transaction représentative de la saisie, un call sur moyenne arithmétique portant sur 64 fixings d'un \omterm{def:qm:stochastic-differential-equations:gbm}{mouvement brownien géométrique} ($S_0 = K = 100$, $r = 3\%$, $\sigma = 25\%$, un an), $2^{14}$ trajectoires donnent
6.495 avec une \omterm{def:qm:estimation:estimator}{erreur type} de 0.077~; le prix vaut 6.4788.

\begin{definition}[Générateurs pseudo-aléatoires et à compteur, méthode d'inversion]\label{def:qm:monte-carlo:prng}
Un \emph{générateur de nombres pseudo-aléatoires}\index{générateur de nombres pseudo-aléatoires} est une récurrence déterministe $s_{k+1} = F(s_k)$, $u_k = G(s_k)$ dont les sorties passent les tests statistiques
d'uniformité indépendante~; le Mersenne Twister (Matsumoto et Nishimura, 1998) et PCG (O'Neill, 2014~; générateur par défaut de NumPy) en sont des exemples. Un \emph{générateur à compteur}\index{générateur à compteur}
calcule $u = G_k(c)$ directement à partir d'un compteur $c$ et d'une clé $k$ grâce à une bijection $G_k$, de sorte que tout tirage peut être produit sans ses prédécesseurs~; Philox (Salmon, Moraes, Dror
et Shaw, 2011) en est un. La \emph{méthode d'inversion}\index{méthode d'inversion} transforme une variable uniforme $U$ en un tirage $F^{-1}(U)$ de la fonction de répartition $F$.
\end{definition}

Un \omterm{def:qm:monte-carlo:prng}{générateur à compteur} rend un lot distribué reproductible. Le moteur attribue à la trajectoire $j$ de la transaction $i$ le compteur $(b, j, i, 0)$ pour son $b$-ième bloc de quatre nombres, de sorte que
les nombres de la trajectoire ne dépendent ni de la répartition des trajectoires entre les machines, ni de l'ordre dans lequel elles s'exécutent, ni de leur nombre~: relancer une seule trajectoire pour déboguer une transaction la reproduit à l'identique, bit pour
bit. Le Philox4x64-10 du kit reproduit l'implémentation de NumPy bloc par bloc, en Python, en C++20 et en Rust (la version C++ figure dans la \cref{tut:qm:monte-carlo:asian}). Les normales proviennent de la \omterm{def:qm:monte-carlo:prng}{méthode d'inversion} avec
l'approximation rationnelle d'Acklam de $\Phi^{-1}$ (erreur relative inférieure à $1.15 \times 10^{-9}$), qui, contrairement à Box--Muller, associe une uniforme à une normale~: la propriété dont ont besoin les points à discrépance faible.

\section{Réduction de variance et échantillonnage préférentiel}

\begin{definition}[Variables antithétiques, variable de contrôle, échantillonnage stratifié]\label{def:qm:monte-carlo:vr}
Les \emph{variables antithétiques}\index{variables antithétiques} associent chaque tirage $Z$ à $-Z$ (chaque uniforme $U$ à $1 - U$) et moyennent $\frac12(f(Z) + f(-Z))$. Une \emph{variable de contrôle}\index{variable de contrôle}
est une variable $C$ dont l'espérance est connue, utilisée dans $\hat\mu_\beta = \frac1n\sum_i(Y_i - \beta(C_i - \E C))$. L'\emph{échantillonnage stratifié}\index{echantillonnage@échantillonnage stratifié} partitionne l'espace d'une entrée en
strates de probabilité connue $p_s$ et tire un nombre fixe $n_s = np_s$ dans chacune.
\end{definition}

\begin{proposition}[Coefficient optimal de la variable de contrôle]\label{prop:qm:monte-carlo:cv}
$\Var(Y - \beta(C - \E C))$ est minimale en $\beta^* = \Cov(Y, C)/\Var(C)$, où elle vaut $(1 - \rho^2)\Var(Y)$ avec $\rho = \mathrm{Corr}(Y, C)$.
\end{proposition}

\begin{proof}
$\Var(Y - \beta C) = \Var Y - 2\beta\Cov(Y, C) + \beta^2\Var C$ est un trinôme en $\beta$ minimal en $\beta^*$~; en substituant, on obtient $\Var Y - \Cov(Y, C)^2/\Var C = (1 - \rho^2)\Var Y$.
\end{proof}

Pour le call sur moyenne, la moyenne géométrique $G = (\prod_kS_{t_k})^{1/64}$ suit une loi lognormale et son call admet une formule fermée (Kemna et Vorst, 1990)~: 6.1793. Son gain a une corrélation de 0.99941 avec
celui de la moyenne arithmétique, donc $1 - \rho^2 = 0.0012$ et, avec $\beta = 1.04$ estimé sur les mêmes trajectoires, la variance est divisée par 853 à nombre de trajectoires égal. Les paires antithétiques font beaucoup moins bien~:
les deux gains d'une paire ont une corrélation de $-0.43$, soit un facteur de 1.8 à nombre d'évaluations de gain égal. La stratification est l'idée de la section suivante~: elle élimine la variance entre strates, et un
ensemble de points à discrépance faible stratifie chaque coordonnée, ainsi que beaucoup de leurs produits, d'un seul coup.

\begin{definition}[Échantillonnage préférentiel]\label{def:qm:monte-carlo:is}
L'\emph{échantillonnage préférentiel}\index{echantillonnage@échantillonnage préférentiel} tire $X$ selon une densité $g$ au lieu de $f$ et moyenne $h(X)f(X)/g(X)$, c'est-à-dire le gain multiplié par le rapport de vraisemblance (une \omterm{def:qm:probability-at-speed:change}{dérivée de Radon--Nikodym})~;
il est sans biais dès que $g > 0$ là où $hf \ne 0$.
\end{definition}

\begin{proposition}[Densité d'importance à variance nulle]\label{prop:qm:monte-carlo:zero}
Pour $h \ge 0$ avec $\mu = \E_f[h(X)] > 0$, la densité $g^* = hf/\mu$ vérifie $h(X)f(X)/g^*(X) = \mu$ pour chaque tirage~: l'\omterm{def:qm:estimation:estimator}{estimateur} est de variance nulle.
\end{proposition}

\begin{proof}
Il suffit de substituer~: $hf/g^* = hf\mu/(hf) = \mu$ partout où $h > 0$, et $g^*$ ne porte aucune masse là où $h = 0$.
\end{proof}

La densité optimale exige la réponse, mais elle indique où chercher~: il faut échantillonner proportionnellement à l'endroit où vit le gain. Pour un digital qui paie si $S_T > 250$ (probabilité $1.21 \times 10^{-4}$, le
driver normal au-delà de $z^* = 3.67$), 100\,000 tirages ordinaires produisent six succès et une \omterm{def:qm:estimation:estimator}{erreur type} relative de 20~\%. En décalant la moyenne du driver de $\theta$ et en pondérant chaque tirage par $e^{-\theta Z +
\theta^2/2}$, on ramène cette erreur à 0.64~\% pour $\theta = 3.8$, soit un facteur de variance d'environ 1\,000 (\cref{fig:qm:monte-carlo:is})~; un décalage excessif la dégrade de nouveau, car les poids des rares tirages qui
comptent explosent.

\begin{omfigure}
\begin{tikzpicture}
\begin{semilogyaxis}[width=10cm,height=4.8cm,xlabel={décalage de moyenne $\theta$ du driver normal},ylabel={erreur type relative},
  tick label style={font=\small},label style={font=\small},xmin=0,xmax=5,ymin=3e-3,ymax=0.5,grid=major,grid style={gray!25},table/col sep=comma]
\addplot[omDef,very thick,mark=*,mark size=1.4pt] table[x=theta,y=relse]{figdata/methods/26-monte-carlo/is.csv};
\draw[omThm,dashed,thick] (axis cs:3.67,3e-3) -- (axis cs:3.67,0.5) node[pos=0.85,right,font=\footnotesize]{$z^* = 3.67$};
\end{semilogyaxis}
\end{tikzpicture}

\omcaption{\omterm{def:qm:monte-carlo:is}{Échantillonnage préférentiel} pour un digital qui paie si $S_T > 250$ ($S_0 = 100$, $\sigma = 25\%$, un an~; probabilité $1.21 \times 10^{-4}$)~: \omterm{def:qm:estimation:estimator}{erreur type} relative de 100\,000 tirages en fonction du
décalage de moyenne du driver normal. Données~: le tutoriel du chapitre, avec graine fixée.}\label{fig:qm:monte-carlo:is}
\end{omfigure}

\section{Suites à discrépance faible}

\begin{definition}[Suite à discrépance faible, suite de Sobol, quasi-Monte-Carlo]\label{def:qm:monte-carlo:qmc}
La discrépance étoilée $D_n^*$ de points $x_1, \dots, x_n$ de $[0, 1)^d$ est la plus grande différence, sur les pavés $[0, a)$, entre la fraction des points situés dans le pavé et son volume. Une \emph{suite à
discrépance faible}\index{suite à discrépance faible} vérifie $D_n^* = O((\log n)^d/n)$. La \emph{suite de Sobol}\index{suite de Sobol} (Sobol, 1967) est une suite à discrépance faible en base 2, construite coordonnée par
coordonnée à partir de polynômes primitifs et de nombres de direction, par ou exclusif sur les chiffres binaires. Le \emph{quasi-Monte-Carlo}\index{quasi-Monte-Carlo} remplace les tirages aléatoires de Monte-Carlo par de tels
points.
\end{definition}

\begin{proposition}[Inégalité de Koksma--Hlawka]\label{prop:qm:monte-carlo:kh}
Pour $f$ à variation bornée $V(f)$ au sens de Hardy et Krause, $\bigl|\frac1n\sum_if(x_i) - \int_{[0,1)^d}f\bigr| \le V(f)\,D_n^*$.
\end{proposition}

La borne donne au \omterm{def:qm:monte-carlo:qmc}{quasi-Monte-Carlo} un taux voisin de $1/n$ pour des intégrandes régulières de dimension effective modérée, mais c'est une borne dans le pire des cas, qui ne peut pas être estimée à partir des points. La randomisation
restaure une estimation de l'erreur.

\begin{definition}[Quasi-Monte-Carlo randomisé, construction par pont brownien]\label{def:qm:monte-carlo:rqmc}
Le \emph{quasi-Monte-Carlo randomisé}\index{quasi-Monte-Carlo randomisé} randomise un ensemble à discrépance faible de sorte que chaque point soit uniforme tandis que l'ensemble conserve sa structure, et estime l'erreur
à partir de réplications indépendantes. Le kit utilise le brouillage uniforme emboîté d'Owen (1995)~: chaque chiffre binaire de chaque coordonnée est retourné par un bit aléatoire attaché au nœud de l'arbre binaire que définissent ses
chiffres de tête. La \emph{construction par pont brownien}\index{construction par pont brownien} affecte la première coordonnée à $W_T$, la suivante à $W_{T/2}$ conditionnellement à ses extrémités, et ainsi de suite
aux niveaux inférieurs, de sorte que les premières coordonnées, là où les points à discrépance faible sont les meilleurs, portent l'essentiel de la variance de la trajectoire.
\end{definition}

La \Cref{fig:qm:monte-carlo:points} montre pourquoi~: 256 points de Sobol brouillés en deux dimensions placent exactement un point dans chacune des 256 boîtes dyadiques de toute forme, alors que des tirages indépendants laissent des trous
et des amas. Sur le call à prix moyen avec $2^{14}$ trajectoires et 32 brouillages indépendants, les points de Sobol brouillés, avec les accroissements dans l'ordre temporel, réduisent la variance d'un facteur d'environ 34~; avec
la \omterm{def:qm:monte-carlo:rqmc}{construction par pont brownien}, d'environ 2\,500~; avec le pont et la \omterm{def:qm:monte-carlo:vr}{variable de contrôle} ensemble, d'environ 37\,000, soit une \omterm{def:qm:estimation:estimator}{erreur type} 190 fois plus petite pour le même nombre de trajectoires. L'erreur
décroît aussi plus vite que $n^{-1/2}$~: les pentes ajustées sont comprises entre $-0.7$ et $-0.8$ de $2^6$ à $2^{15}$ trajectoires (\cref{fig:qm:monte-carlo:cost}).

\begin{omfigure}
\begin{tikzpicture}
\begin{axis}[name=l,width=5.6cm,height=5.6cm,title={Philox},title style={font=\small},xmin=0,xmax=1,ymin=0,ymax=1,
  tick label style={font=\scriptsize},xtick={0,0.5,1},ytick={0,0.5,1},table/col sep=comma]
\addplot[only marks,mark=*,mark size=0.9pt,omThm] table[x=px,y=py]{figdata/methods/26-monte-carlo/points.csv};
\end{axis}
\begin{axis}[at={(l.east)},anchor=west,xshift=1cm,width=5.6cm,height=5.6cm,title={Sobol brouillé},title style={font=\small},xmin=0,xmax=1,ymin=0,ymax=1,
  tick label style={font=\scriptsize},xtick={0,0.5,1},ytick={0,0.5,1},table/col sep=comma]
\addplot[only marks,mark=*,mark size=0.9pt,omDef] table[x=sx,y=sy]{figdata/methods/26-monte-carlo/points.csv};
\end{axis}
\end{tikzpicture}

\omcaption{256 points dans le carré unité~: uniformes Philox indépendantes et deux premières coordonnées de points de Sobol brouillés par Owen. Données~: le tutoriel du chapitre,
avec graine fixée.}\label{fig:qm:monte-carlo:points}
\end{omfigure}

\begin{omfigure}
\begin{tikzpicture}
\begin{loglogaxis}[width=11cm,height=6cm,xlabel={nombre de trajectoires},ylabel={erreur type du prix},
  tick label style={font=\small},label style={font=\small},xmin=50,xmax=45000,ymin=1e-4,ymax=3,grid=major,grid style={gray!25},
  legend columns=3,legend cell align=left,legend style={font=\footnotesize,at={(0.5,-0.25)},anchor=north,draw=none,fill=none,
  /tikz/every even column/.append style={column sep=6pt}},table/col sep=comma]
\addplot[omThm,very thick,mark=*,mark size=1.4pt] table[x=n,y=plain]{figdata/methods/26-monte-carlo/cost.csv};
\addplot[omThm!50,very thick,mark=o,mark size=1.6pt] table[x=n,y=rqmcinc]{figdata/methods/26-monte-carlo/cost.csv};
\addplot[black,very thick,mark=square*,mark size=1.4pt] table[x=n,y=cv]{figdata/methods/26-monte-carlo/cost.csv};
\addplot[omDef!60,very thick,mark=triangle*,mark size=1.8pt] table[x=n,y=rqmcbridge]{figdata/methods/26-monte-carlo/cost.csv};
\addplot[omDef,very thick,mark=diamond*,mark size=1.8pt] table[x=n,y=rqmccv]{figdata/methods/26-monte-carlo/cost.csv};
\legend{Monte Carlo classique,{RQMC, ordre temporel},variable de contrôle,{RQMC, pont},{RQMC, pont + contrôle}}
\end{loglogaxis}
\end{tikzpicture}

\omcaption{\omterm{def:qm:estimation:estimator}{Erreur type} du call à prix moyen (64 constatations) en fonction du nombre de trajectoires~: Monte Carlo classique et \omterm{def:qm:monte-carlo:vr}{variable de contrôle} (pente $-\frac12$), et points de Sobol brouillés (16
brouillages indépendants) avec les accroissements dans l'ordre temporel, avec la \omterm{def:qm:monte-carlo:rqmc}{construction par pont brownien}, et avec le pont et la \omterm{def:qm:monte-carlo:vr}{variable de contrôle}. Données~: le tutoriel du chapitre,
avec graine fixée.}\label{fig:qm:monte-carlo:cost}
\end{omfigure}

\section{Discrétisation et méthodes multiniveaux}

\begin{definition}[Ordres de convergence forts et faibles, schéma de Milstein]\label{def:qm:monte-carlo:order}
Un schéma de pas $h$ possède un \emph{ordre de convergence fort}\index{ordre de convergence fort} $\gamma$ si $\E|X_T^h - X_T| \le Ch^\gamma$, et un \emph{ordre de convergence faible}\index{ordre de convergence faible} $\beta$ si
$|\E g(X_T^h) - \E g(X_T)| \le Ch^\beta$ pour $g$ régulière. Le \emph{schéma de Milstein}\index{schéma de Milstein} pour $dX = a(X)\,dt + b(X)\,dW$ ajoute au pas d'Euler--Maruyama le
terme $\frac12b(X)b'(X)((\Delta W)^2 - h)$.
\end{definition}

\begin{proposition}[Ordres d'Euler et de Milstein]\label{prop:qm:monte-carlo:orders}
Sous des conditions de Lipschitz et de croissance (et la régularité de $a$, $b$ pour Milstein), le \omterm{def:qm:stochastic-differential-equations:euler}{schéma d'Euler--Maruyama} est d'ordre fort $\frac12$ et d'ordre faible 1, et le \omterm{def:qm:monte-carlo:order}{schéma de Milstein} d'ordre fort 1
(Kloeden et Platen, 1992).
\end{proposition}

\begin{proof}[Idée]
Le pas d'Euler néglige l'intégrale stochastique double $\int\!\!\int dW\,dW = \frac12((\Delta W)^2 - h)$, de taille $h$ et d'espérance nulle, multipliée par $bb'$~; ces erreurs s'additionnent comme une marche aléatoire sur $1/h$
pas, soit $h\sqrt{1/h} = h^{1/2}$ trajectoire par trajectoire, tandis que leurs espérances se compensent, laissant un biais en loi d'ordre $O(h)$. Milstein conserve ce terme, et les termes manquants suivants sont d'ordre $h^{3/2}$ par pas.
\end{proof}

Un prix exige l'ordre faible~; une couverture simulée le long de la trajectoire, ou un \omterm{def:qm:estimation:estimator}{estimateur} multiniveau, exige l'ordre fort. Sur un \omterm{def:qm:stochastic-differential-equations:gbm}{mouvement brownien géométrique} avec $\sigma = 50\%$, piloté par les mêmes accroissements que
la solution exacte, les pentes fortes ajustées valent 0.49 pour Euler et 0.89 pour Milstein, et l'erreur faible de la moyenne vaut exactement $X_0|(1 + \mu h)^n - e^{\mu T}|$, pente 1.00
(\cref{fig:qm:monte-carlo:orders}).

\begin{omfigure}
\begin{tikzpicture}
\begin{loglogaxis}[width=10cm,height=5cm,xlabel={pas de temps $h$ (années)},ylabel={erreur},x dir=reverse,
  tick label style={font=\small},label style={font=\small},xmin=0.003,xmax=0.6,ymin=3e-4,ymax=20,grid=major,grid style={gray!25},
  legend columns=3,legend cell align=left,legend style={font=\footnotesize,at={(0.5,-0.3)},anchor=north,draw=none,fill=none,
  /tikz/every even column/.append style={column sep=8pt}},table/col sep=comma]
\addplot[omThm,very thick,mark=*,mark size=1.4pt] table[x=h,y=euler]{figdata/methods/26-monte-carlo/orders.csv};
\addplot[omDef,very thick,mark=square*,mark size=1.4pt] table[x=h,y=milstein]{figdata/methods/26-monte-carlo/orders.csv};
\addplot[black,very thick,mark=triangle*,mark size=1.8pt] table[x=h,y=weak]{figdata/methods/26-monte-carlo/orders.csv};
\legend{erreur forte (Euler),erreur forte (Milstein),erreur faible de la moyenne}
\end{loglogaxis}
\end{tikzpicture}

\omcaption{Erreur forte $\E|X_T^h - X_T|$ des schémas d'Euler et de Milstein (20\,000 trajectoires d'un \omterm{def:qm:stochastic-differential-equations:gbm}{mouvement brownien géométrique}, $\mu = 5\%$, $\sigma = 50\%$, $X_0 = 100$, un an) et erreur faible exacte de
la moyenne, en fonction du pas de temps. Données~: le tutoriel du chapitre, avec graine fixée.}\label{fig:qm:monte-carlo:orders}
\end{omfigure}

\begin{definition}[Monte-Carlo multiniveau]\label{def:qm:monte-carlo:mlmc}
Le \emph{Monte-Carlo multiniveau}\index{Monte-Carlo multiniveau} (Giles, 2008) écrit l'espérance au niveau le plus fin comme une somme télescopique $\E P_L = \E P_0 + \sum_{l=1}^L\E[P_l - P_{l-1}]$, où $P_l$ est le
gain calculé avec le pas $h_l = 2^{-l}T$, et estime chaque terme indépendamment, $P_l$ et $P_{l-1}$ dans chaque correction étant calculés à partir de la même trajectoire brownienne.
\end{definition}

\begin{proposition}[Complexité du Monte-Carlo multiniveau]\label{prop:qm:monte-carlo:giles}
Si le biais de $P_l$ est $O(2^{-\alpha l})$, la variance de la correction $V_l = O(2^{-\beta l})$ et son coût $C_l = O(2^{\gamma l})$, avec $\alpha \ge \frac12\min(\beta, \gamma)$, alors une erreur quadratique moyenne
$\epsilon^2$ coûte $O(\epsilon^{-2})$ si $\beta > \gamma$, $O(\epsilon^{-2}(\log\epsilon)^2)$ si $\beta = \gamma$ et $O(\epsilon^{-2 - (\gamma - \beta)/\alpha})$ si $\beta < \gamma$~; le Monte-Carlo standard coûte
$O(\epsilon^{-2 - \gamma/\alpha})$.
\end{proposition}

\begin{proof}[Idée]
Minimiser le coût $\sum_lN_lC_l$ sous la contrainte de variance $\sum_lV_l/N_l = \epsilon^2/2$ donne $N_l \propto\sqrt{V_l/C_l}$ et un coût total
\[2\epsilon^{-2}\Bigl(\sum_l\sqrt{V_lC_l}\Bigr)^2;\]
la somme est dominée par les niveaux les plus grossiers
si $\beta > \gamma$, par les plus fins si $\beta < \gamma$, et croît comme $L \sim \log\epsilon^{-1}$ si elles sont égales. Le niveau le plus fin $L$ est fixé par le biais, $2^{-\alpha L} \sim \epsilon$.
\end{proof}

L'ordre fort fixe $\beta$~: Milstein sur un gain lipschitzien donne $\beta = 2 > \gamma = 1$, et les variances des corrections dans le tutoriel diminuent d'un facteur d'environ quatre par niveau
(0.23, 0.060, 0.016, 0.0045, 0.0012, 0.00030, 0.000079). Pour un call européen ($\alpha = 1$), le pilote listé à la \cref{tut:qm:monte-carlo:asian} atteint une erreur quadratique moyenne de 0.005 avec 7 niveaux
pour un coût de $3.2 \times 10^7$ pas de temps, là où le Monte-Carlo standard au même pas le plus fin en exigerait $3.0 \times 10^9$~: un facteur de 94, qui croît à mesure que $\epsilon$ diminue (8, 20, 45, 94 pour
$\epsilon$ de 0.04 à 0.005~; \cref{fig:qm:monte-carlo:mlmc}).

\begin{omfigure}
\begin{tikzpicture}
\begin{loglogaxis}[width=10cm,height=5cm,xlabel={erreur quadratique moyenne cible $\epsilon$},ylabel={coût (pas de temps)},x dir=reverse,
  tick label style={font=\small},label style={font=\small},xmin=0.004,xmax=0.05,ymin=1e5,ymax=1e10,grid=major,grid style={gray!25},
  xtick={0.005,0.01,0.02,0.04},xticklabels={0.005,0.01,0.02,0.04},
  legend columns=2,legend cell align=left,legend style={font=\footnotesize,at={(0.5,-0.3)},anchor=north,draw=none,fill=none,
  /tikz/every even column/.append style={column sep=8pt}},table/col sep=comma]
\addplot[omThm,very thick,mark=*,mark size=1.6pt] table[x=eps,y=standard]{figdata/methods/26-monte-carlo/mlmc.csv};
\addplot[omDef,very thick,mark=square*,mark size=1.6pt] table[x=eps,y=mlmc]{figdata/methods/26-monte-carlo/mlmc.csv};
\legend{Monte Carlo standard (Milstein),Monte-Carlo multiniveau}
\end{loglogaxis}
\end{tikzpicture}

\omcaption{Coût de la valorisation d'un call européen ($S_0 = K = 100$, $r = 3\%$, $\sigma = 25\%$, un an) pour une erreur quadratique moyenne cible~: \omterm{def:qm:monte-carlo:mlmc}{Monte-Carlo multiniveau} avec corrections de Milstein, contre
Monte Carlo standard au même pas le plus fin. Données~: le tutoriel du chapitre, avec graine fixée.}\label{fig:qm:monte-carlo:mlmc}
\end{omfigure}

\section{Tutoriel~: seize fois moins de machines}\label{tut:qm:monte-carlo:asian}

\begin{tutorial}
\textbf{Objectif.} Valoriser la transaction représentative du lot de cinq façons et mesurer la variance de chaque \omterm{def:qm:estimation:estimator}{estimateur} à nombre de trajectoires égal. \textbf{État final~:} les facteurs 1.8, 853, 34, 2\,500 et 37\,000 et
la \cref{fig:qm:monte-carlo:cost}.
\begin{tutsteps}
\item \textbf{Variables uniformes à compteur.} Philox4x64-10 en C++20, reproduisant le bloc de \texttt{Philox} de NumPy bloc par bloc.
\omcode{code/firm/mcengine/cpp/firm_mcengine.hpp}{113}{124}{Philox4x64-10~: dix tours de deux multiplications larges (C++20).}\label{lst:qm:monte-carlo:philox}
\item \textbf{Points de Sobol brouillés.} Points de Sobol en code de Gray à partir des nombres de direction de Joe et Kuo, et brouillage d'Owen par nœuds d'arbre hachés.
\omcode{code/firm/mcengine/firm_mcengine.py}{268}{297}{Points de Sobol et brouillage uniforme imbriqué (Python).}\label{lst:qm:monte-carlo:sobol}
\item \textbf{\omterm{def:qm:estimation:estimator}{Estimateurs}.} Exécuter \texttt{table()} dans \texttt{qm\_mc.py}~: \omterm{def:qm:estimation:estimator}{estimateur} standard, \omterm{def:qm:monte-carlo:vr}{variables antithétiques}, \omterm{def:qm:monte-carlo:vr}{variable de contrôle}, et points de Sobol brouillés avec et sans le pont (32 brouillages de $2^{14}$ points)~; puis
\texttt{error\_vs\_cost()}.
\item \textbf{Multiniveau.} Exécuter \texttt{mlmc\_study()}~; le pilote figure ci-dessous.
\omcode{code/firm/mcengine/firm_mcengine.py}{384}{412}{Pilote de Monte-Carlo multiniveau de Giles (Python).}\label{lst:qm:monte-carlo:mlmc}
\end{tutsteps}
\textbf{Pour aller plus loin.} Fixer le strike à 130 et observer la baisse du facteur de la \omterm{def:qm:monte-carlo:vr}{variable de contrôle} avec la corrélation~; remplacer le pont par une construction en composantes principales~; valoriser une option à barrière avec
la probabilité de franchissement du pont du chapitre~2 à l'intérieur du pilote multiniveau.
\end{tutorial}

\section{Construction~: le moteur de Monte-Carlo, étape trois}\label{bld:qm:monte-carlo:mcengine}

\begin{build}
\bfield{Rôle} Rendre chaque trajectoire plus utile et chaque trajectoire reproductible~: les générateurs, ensembles de points et \omterm{def:qm:estimation:estimator}{estimateurs} qu'appellent les lots de valorisation et de risque de la firme.
\bfield{Interface} \texttt{philox4x64}, \texttt{philox\_uniforms(n, dim, seed, stream, first\_path)}~; \texttt{sobol\_points}, \texttt{owen\_scramble}~; \texttt{norm\_ppf}~;
\texttt{bridge\_from\_normals}~; \texttt{asian\_payoffs}, \texttt{geometric\_asian\_price}, \texttt{control\_variate}~; \texttt{mlmc(sampler, eps)}. Jumeaux C++20 et Rust de \texttt{philox4x64},
\texttt{sobol\_point} (huit dimensions) et \texttt{owen\_scramble}.
\bfield{Règles} Les compteurs, et non un état séquentiel, identifient les trajectoires~; chaque estimation est accompagnée de son \omterm{def:qm:estimation:estimator}{erreur type} (calculée à partir de réplications pour le \omterm{def:qm:monte-carlo:rqmc}{quasi-Monte-Carlo randomisé}, jamais à l'intérieur d'un seul ensemble de points)~;
les points de Sobol non brouillés ne servent jamais aux barres d'erreur~; les nombres de direction portent leur licence.
\bfield{Tests d'acceptation} \texttt{code/firm/mcengine/\{tests,cpp,rust\}}~: Philox face à celui de NumPy~; la découpe des trajectoires les reproduit~; les premiers points publiés de Joe et Kuo~; des entiers brouillés identiques dans
les trois langages~; la propriété de réseau après brouillage~; $\Phi^{-1}$ face à la bibliothèque standard à $10^{-8}$ près~; la covariance du pont~; la forme fermée géométrique face à la simulation et face à
Black--Scholes pour une seule date de fixing~; le pilote multiniveau sur un échantillonneur à réponse connue.
\bfield{Pour aller plus loin} Les 21\,201 dimensions complètes de Joe et Kuo~; une construction de trajectoires en composantes principales~; le Monte Carlo conditionnel pour les barrières~; les grecques adjointes le long des trajectoires (chapitre~28).
\end{build}

\begin{omsources}
\item N.~Metropolis et S.~Ulam, « The Monte Carlo method », \emph{Journal of the American Statistical Association} 44, 1949~; P.~P.~Boyle, « Options: a Monte Carlo approach », \emph{Journal of
Financial Economics} 4, 1977.
\item P.~Glasserman, \emph{Monte Carlo Methods in Financial Engineering}, Springer, 2003.
\item J.~K.~Salmon, M.~A.~Moraes, R.~O.~Dror et D.~E.~Shaw, « Parallel random numbers: as easy as 1, 2, 3 », SC11, 2011~; M.~Matsumoto et T.~Nishimura, \emph{ACM TOMACS} 8, 1998~; M.~E.~O'Neill,
PCG, Harvey Mudd College report HMC-CS-2014-0905, 2014.
\item I.~M.~Sobol', \emph{USSR Computational Mathematics and Mathematical Physics} 7, 1967~; S.~Joe et F.~Y.~Kuo, \emph{SIAM Journal on Scientific Computing} 30, 2008 (nombres de direction
\texttt{new-joe-kuo-6.21201})~; A.~B.~Owen, « Randomly permuted $(t, m, s)$-nets and $(t, s)$-sequences », 1995~; B.~Moskowitz et R.~E.~Caflisch, \emph{Mathematical and Computer Modelling} 23, 1996.
\item A.~G.~Z.~Kemna et A.~C.~F.~Vorst, \emph{Journal of Banking and Finance} 14, 1990.
\item M.~B.~Giles, « Multilevel Monte Carlo path simulation », \emph{Operations Research} 56, 2008~; P.~E.~Kloeden et E.~Platen, \emph{Numerical Solution of Stochastic Differential Equations},
Springer, 1992~; G.~N.~Milstein, \emph{Theory of Probability and its Applications} 19, 1975.
\item P.~J.~Acklam, « An algorithm for computing the inverse normal cumulative distribution function » (note en ligne, 2004).
\end{omsources}

\section{Exercices}

\begin{exercise}[$\star$]\label{exo:qm:monte-carlo:1}
Un lot a une \omterm{def:qm:estimation:estimator}{erreur type} de 80\,000 dollars avec 100\,000 trajectoires. Combien de trajectoires permettent d'atteindre 20\,000, et quel facteur de \omterm{def:qm:monte-carlo:mc}{réduction de variance} le même lot à 100\,000 trajectoires donne-t-il~?
\end{exercise}

\begin{exercise}[$\star$]\label{exo:qm:monte-carlo:2}
Une \omterm{def:qm:monte-carlo:vr}{variable de contrôle} a une corrélation de 0.9 avec le payoff. Quel est le facteur de variance~? Et avec 0.99~?
\end{exercise}

\begin{exercise}[$\star$]\label{exo:qm:monte-carlo:3}
Pourquoi les \omterm{def:qm:monte-carlo:vr}{variables antithétiques} aident-elles un payoff monotone et non une fonction paire de $Z$, comme $Z^2$~?
\end{exercise}

\begin{exercise}[$\star\star$]\label{exo:qm:monte-carlo:4}
Montrer que la stratification proportionnelle n'augmente jamais la variance~: avec des strates $s$ de probabilité $p_s$ et des variances conditionnelles $\sigma_s^2$, la variance stratifiée vaut $\frac1n\sum_sp_s\sigma_s^2 \le
\sigma^2/n$.
\end{exercise}

\begin{exercise}[$\star\star$]\label{exo:qm:monte-carlo:5}
Pour $P(Z > z^*)$ avec $Z$ normale centrée réduite, calculer le second moment de l'\omterm{def:qm:estimation:estimator}{estimateur} par \omterm{def:qm:monte-carlo:is}{échantillonnage préférentiel} avec le décalage $\theta$, et montrer qu'il est minimisé près de $\theta = z^*$ pour de grands $z^*$.
\end{exercise}

\begin{exercise}[$\star\star$]\label{exo:qm:monte-carlo:6}
Dériver le pas de Milstein pour le \omterm{def:qm:stochastic-differential-equations:gbm}{mouvement brownien géométrique} et montrer que c'est le développement au second ordre du pas exact $X e^{(\mu - \sigma^2/2)h + \sigma\Delta W}$ en $\Delta W$.
\end{exercise}

\begin{exercise}[$\star\star\star$]\label{exo:qm:monte-carlo:7}
\emph{Codage.} Calculer le facteur de variance de la \omterm{def:qm:monte-carlo:vr}{variable de contrôle} pour les strikes 80, 100 et 130, et expliquer la tendance.
\end{exercise}

\begin{exercise}[$\star\star\star$]\label{exo:qm:monte-carlo:8}
\emph{Repérer l'erreur.} « Nous avons basculé le lot vers des points de Sobol, 100\,000 par transaction, et nous rapportons l'\omterm{def:qm:estimation:estimator}{erreur type} $s/\sqrt n$ calculée à partir des points comme auparavant~; elle a été divisée par dix. »
\end{exercise}

\section{Problème~: seize fois moins de machines}

\begin{problem}[{Problème du week-end --- un facteur quatre sur l'erreur, sans seize fois plus de trajectoires}]\label{pb:qm:monte-carlo:1}
La transaction représentative est un call sur moyenne arithmétique sur 64 fixings ($S_0 = K = 100$, $r = 3\%$, $\sigma = 25\%$, un an), valorisé avec $2^{14}$ trajectoires~; le lot a besoin d'un facteur quatre sur l'\omterm{def:qm:estimation:estimator}{erreur type}.

\medskip\noindent\textbf{Partie I --- L'\omterm{def:qm:estimation:estimator}{estimateur} standard.}
\begin{enumerate}
\item Quels sont l'estimation standard, son \omterm{def:qm:estimation:estimator}{erreur type} et le prix de référence~?
\item Quel facteur de variance la cible exige-t-elle~?
\item Que dit le produit coût-variance pour comparer des \omterm{def:qm:estimation:estimator}{estimateurs}~?
\item Pourquoi le \omterm{def:qm:monte-carlo:prng}{générateur à compteur} rend-il le lot reproductible d'une machine à l'autre~?
\item Pourquoi utiliser l'inversion plutôt que Box--Muller pour alimenter les points à discrépance faible~?
\end{enumerate}

\medskip\noindent\textbf{Partie II --- \omterm{def:qm:monte-carlo:mc}{Réduction de variance}.}
\begin{enumerate}[resume]
\item Quel facteur donnent les paires antithétiques, et pourquoi si peu~?
\item Qu'est-ce que la \omterm{def:qm:monte-carlo:vr}{variable de contrôle}, quelle est sa valeur fermée, sa corrélation et son facteur~?
\item Quel est le coefficient optimal, et que cela a-t-il estimé~?
\item Où l'\omterm{def:qm:monte-carlo:is}{échantillonnage préférentiel} aide-t-il, et de combien sur le digital profondément hors de la monnaie~?
\item Pourquoi un décalage trop grand l'empire-t-il~?
\end{enumerate}

\medskip\noindent\textbf{Partie III --- Points et trajectoires.}
\begin{enumerate}[resume]
\item Quel facteur donnent les points de Sobol mélangés avec les accroissements dans l'ordre temporel, et avec le pont~?
\item Pourquoi le pont compte-t-il autant~?
\item Que donnent les deux ensemble, et quelle est la pente de l'erreur en fonction du nombre de trajectoires~?
\item Comment estime-t-on l'erreur du \omterm{def:qm:monte-carlo:rqmc}{quasi-Monte-Carlo randomisé}~?
\item Que mesurent les ordres forts et faibles, et quelles sont leurs valeurs ajustées ici~?
\end{enumerate}

\medskip\noindent\textbf{Partie IV --- Jugement.}
\begin{enumerate}[resume]
\item Que économise le \omterm{def:qm:monte-carlo:mlmc}{Monte-Carlo multiniveau} sur le call européen à $\epsilon = 0.005$, et pourquoi~?
\item Le facteur de 37\,000 se maintiendrait-il sur l'ensemble du lot~?
\item Que doit consigner le lot pour que les trajectoires d'une transaction puissent être régénérées~?
\item Énoncer le \emph{résultat nommé}~: le facteur de \omterm{def:qm:monte-carlo:mc}{réduction de variance} de la \omterm{def:qm:monte-carlo:vr}{variable de contrôle} combinée au \omterm{def:qm:monte-carlo:rqmc}{quasi-Monte-Carlo randomisé} sur la transaction représentative, face aux 16 requis.
\item En une phrase~: où va le travail d'un moteur de Monte-Carlo~?
\end{enumerate}
\end{problem}

\section{Questions d'entretien}

\begin{interviewq}[$\star$ \iqroles{chercheur, développeur}]\label{iq:qm:monte-carlo:1}
Comment l'erreur de Monte-Carlo évolue-t-elle avec le nombre de trajectoires et avec la dimension~?
\end{interviewq}

\begin{interviewq}[$\star\star$ \iqroles{chercheur}]\label{iq:qm:monte-carlo:2}
Citez trois techniques de \omterm{def:qm:monte-carlo:mc}{réduction de variance} et dites quand chacune fonctionne.
\end{interviewq}

\begin{interviewq}[$\star\star$ \iqroles{développeur}]\label{iq:qm:monte-carlo:3}
Comment rendez-vous un lot de Monte-Carlo distribué reproductible quel que soit le nombre de machines~?
\end{interviewq}

\begin{interviewq}[$\star\star$ \iqroles{chercheur}]\label{iq:qm:monte-carlo:4}
Pourquoi le \omterm{def:qm:monte-carlo:qmc}{quasi-Monte-Carlo} peut-il mal se comporter sur une trajectoire à 250 pas, et qu'est-ce qui y remédie~?
\end{interviewq}

\begin{interviewq}[$\star\star$ \iqroles{chercheur, risque}]\label{iq:qm:monte-carlo:5}
Comment estimeriez-vous, par simulation, la probabilité d'une perte au-delà d'un seuil très élevé~?
\end{interviewq}

\begin{interviewq}[$\star\star\star$ \iqroles{chercheur}]\label{iq:qm:monte-carlo:6}
Expliquez le \omterm{def:qm:monte-carlo:mlmc}{Monte-Carlo multiniveau} et pourquoi l'ordre fort du schéma compte pour celui-ci.
\end{interviewq}
